\section{重叠窗口的交换性与方差机制}\label{sec:theory}

本节用可解析的窗口均值说明重叠如何改变派生统计量的联合分布。结论用于解释失校准机制，不给出一般拓扑最大扫描的误报率公式。

\begin{proposition}[局部映射的有限依赖范围]\label{prop:dependency}
设原始观测 $\{X_t\}$ 独立同分布，$Y_i=\Psi\bigl((X_t)_{t\in I_i}\bigr)$，其中 $\Psi$ 为固定可测映射且不使用由整条序列共同估计的参数。则 $\{Y_i\}$ 至多为 $m=\lceil W/s\rceil-1$ 阶依赖。
\end{proposition}

命题~\ref{prop:dependency}是文献\cite{NieYue2026}附录命题~6在原始观测独立同分布条件下的特例；证明见补充材料第~S6.5~节。该上界只界定依赖范围：具体协方差是否为零仍取决于 $\Psi$ 与原始分布，局部依赖也不自动等同于不可交换性。

\begin{proposition}[高斯窗口均值的非交换性]\label{prop:counterexample}
设 $X_t\overset{\mathrm{iid}}{\sim}N(0,\sigma^2)$，$\sigma^2>0$，$1\le s<W$，$M\ge3$，记 $Z_i=W^{-1}\sum_{t\in I_i}X_t$。则对 $h=0,\ldots,M-1$ 及 $i=1,\ldots,M-h$，
\begin{equation}\label{eq:window-cov}
 \Gamma_h=\operatorname{Cov}(Z_i,Z_{i+h})
 =\frac{\sigma^2}{W^2}(W-hs)_+ ,
\end{equation}
且随机向量 $(Z_1,\ldots,Z_M)$ 不具交换性。
\end{proposition}

\begin{proof}
只有被两个窗口共享的观测对协方差有贡献，交集 $I_i\cap I_{i+h}$ 含 $(W-hs)_+$ 个指标，得到式\eqref{eq:window-cov}。又因 $1\le s<W$，
$\Gamma_1=\sigma^2(W-s)/W^2>\sigma^2(W-2s)_+/W^2=\Gamma_2$。交换第 $2$、$3$ 个坐标后，第一坐标与二者的协方差互换，协方差矩阵改变；高斯向量的分布由均值和协方差唯一确定，故交换后的联合分布与原分布不同。
\end{proof}

\begin{remark}
该反例严格限定于窗口均值。常数映射始终产生可交换序列；$M=2$ 时两个等方差高斯窗口均值的协方差矩阵对称，交换坐标不改变分布，故 $M\ge3$ 的条件具有实质作用。
\end{remark}

在同一构造下，令 $\bar Z_M=M^{-1}\sum_iZ_i$、$\Gamma_0=\sigma^2/W$，则 $\operatorname{Var}(\bar Z_M)=\Gamma_0A_{\mathrm{mean},M}/M$，其中
\begin{equation}\label{eq:variance-ratio}
\begin{aligned}
 A_{\mathrm{mean},M}
 &=1+2\sum_{h=1}^{M-1}\Bigl(1-\frac hM\Bigr)\Bigl(1-\frac{hs}{W}\Bigr)_+\\
 &=r-\frac{r^2-1}{3M},
 \qquad r=W/s\in\mathbb N,\quad M\ge r .
\end{aligned}
\end{equation}
在 $W=64$、$s=16$、$M=192$ 时 $r=4$，$A_{\mathrm{mean},M}=4-15/576\approx3.974$。该数值说明原顺序的部分和携带滞后协方差，而随机排列的桥接过程只保留经验边际变异：给定窗口值 $z_1,\ldots,z_M$，记 $s_M^2=M^{-1}\sum_{i=1}^M(z_i-\bar z)^2$，则排列桥接的条件方差为 $k(M-k)s_M^2/[M(M-1)]$，与滞后结构无关（证明见补充材料第~S6.7、S6.8~节）。但 $3.974$ 只刻画标量窗口均值的方差比，不能解释为 BPI 最大扫描 I 类错误率的膨胀倍数——后者还受高维方向、特征谱与极值操作影响。完整非线性管线的校准因此需由第~\ref{sec:h0} 节的模拟评价给出。
